\begin{proposition}[Proved under stated hypotheses: arbitrary boundary and finite local images]
\label{D:finite-local-images-arbitrary}
Let $X$ be a normal connected proper scheme over an algebraically
closed field, let $Y\subset X$ be closed of codimension at least
two, and put $U=X\setminus Y$.  For a geometric point $\bar x$ of
$Y$, put
\[
 U_{\bar x}=U\times_X\operatorname{Spec}
                    \mathcal O^{\mathrm{sh}}_{X,\bar x},
 \qquad
 I_{\bar x}=\operatorname{im}\bigl(\pi_1(U_{\bar x})
                                  \longrightarrow\pi_1(U)\bigr).
\]
If every $I_{\bar x}$ is finite, then $\pi_1(U)$ is topologically
finitely generated.  No uniform bound on the orders is required,
and $Y$ need not be finite.  It is enough to check the hypothesis
at geometric closed points of $Y$.
\end{proposition}
\begin{proof}
We use the following precise local-to-global tower theorem:
for a normal integral excellent scheme $X$, a tower of finite
quasi-\'etale Galois covers which is eventually \'etale after
strict localization at each geometric point is eventually
\'etale globally
\cite[Theorem 3.7, pp.~708--709]{CarvajalRojasStabler2023D}.
Its statement imposes no quasi-projectivity hypothesis.  In our
finite-type setting it also follows from regular alterations and
the finite branch-locus stratification in
\cite[Theorem 2 and \S4, proof of (i)$\Rightarrow$(iii),
pp.~6--7]{Stibitz2017D}.  This is a theorem about all finite
quasi-\'etale covers, with no restriction on their degrees.

Write $G=\pi_1(U)$.  For an open normal subgroup $H\subset G$,
let $V_H\to U$ be its connected finite \'etale cover and let
$X_H\to X$ be the finite normalization in $k(V_H)$.  We claim
that some $H$ satisfies
\[
              \operatorname{FEt}(X_H)
                    \simeq\operatorname{FEt}(V_H).
\]
If this failed for every $H$, we could construct successively
open normal subgroups $G=H_0\supset H_1\supset H_2\supset\cdots$
so that every $X_{H_{n+1}}\to X_{H_n}$ were non-\'etale.  Indeed,
choose a finite \'etale cover of $V_{H_n}$ whose finite
normalization over $X_{H_n}$ is not \'etale, and dominate its
generic field by a finite Galois extension of $k(X)$ arising
from the Galois closure of that cover over $U$.  If the
dominating normalization were \'etale over $X_{H_n}$, its
intermediate normal finite cover would also be \'etale, contrary
to the choice.  Every $X_{H_n}\to X$ is quasi-\'etale because
its branch locus is contained in $Y$.

Fix $\bar x$.  The descending subgroups
$I_{\bar x}\cap H_n$ of the finite group $I_{\bar x}$ stabilize.
The connected components of $V_{H_n}\times_U U_{\bar x}$ are
the covers determined by these intersections.  Once two
successive intersections agree, each component upstairs maps
isomorphically to its corresponding component downstairs.
Finite normalization commutes with \'etale localization, and
normality identifies the normalization of a split generic
algebra with a disjoint union of copies of the normal base.
Consequently $X_{H_{n+1}}\to X_{H_n}$ is \'etale over the
strict localization at $\bar x$ for all sufficiently large
$n$.  At points of $U$ it is already \'etale.  The tower theorem
now contradicts our construction.  This proves the claim.

For this $H$, full faithfulness of restriction follows from
normalization in generic fields, so the displayed equivalence
gives $H\simeq\pi_1(X_H)$.  The scheme $X_H$ is proper;
\cite[X, Theorem 2.9]{SGA1} makes its fundamental group
topologically finitely generated.  Adjoining finitely many
coset representatives generates $G$.

Finally, every point $x\in Y$ specializes to a closed point
$y\in Y$, since $X$ is of finite type over a field.  Choose a
geometric lift of $x$ to
$\operatorname{Spec}\mathcal O^{\mathrm{sh}}_{X,\bar y}$.
Strict localization there is a strict localization of $X$ at
$x$, because the map to $X$ is ind-\'etale.  Thus the map
from $\pi_1(U_{\bar x})$ to $G$ factors through
$\pi_1(U_{\bar y})$, up to the usual base-point conjugacy.
It follows that $I_{\bar x}$ is contained in a conjugate of
$I_{\bar y}$.  Finiteness at geometric closed points suffices.
\end{proof}

\begin{corollary}[Proved under stated hypotheses: klt threefold boundary]
\label{D:klt-threefold-boundary}
Let $k$ be algebraically closed of characteristic $p>5$, let
$X$ be a normal connected proper threefold over $k$, and let
$Y\subset X$ be closed of codimension at least two.  Suppose
that $X$ has klt singularities on a neighborhood of every
point of $Y$.  Then $\pi_1(X\setminus Y)$ is topologically
finitely generated.  The retained singularities are unrestricted,
and neither projectivity nor a local-complete-intersection
hypothesis is required.
\end{corollary}
\begin{proof}
For a geometric closed point $\bar y$ of $Y$, take an affine
neighborhood on which the klt hypothesis holds.  Its strict
localization $S$ and the closed subset $Z=Y\times_X S$ satisfy
Setup 1.1 of \cite[p.~694]{CarvajalRojasStabler2023D}:
$Z$ is arbitrary of codimension at least two, not necessarily
just the closed point or the singular locus.  Their Theorem A
\cite[p.~695; Theorem 4.3 and Corollary 4.5,
pp.~713--714]{CarvajalRojasStabler2023D} states that the full
group $\pi_1(S\setminus Z)$ is finite of order prime to $p$.
Its conclusion is not merely finiteness of the maximal
prime-to-$p$ quotient.  The proof of Corollary 4.5 explicitly
excludes a finite quotient of order divisible by $p$ by passing
to a Sylow $p$-subgroup and applying Corollary 2.15 to the
intermediate klt strict-local cover.  In particular the image
$I_{\bar y}$ in the preceding proposition is finite.
That proposition proves the claim.

This argument uses the published local theorem only on affine
neighborhoods and the separately stated scheme-theoretic tower
theorem globally.  It therefore does not depend on whether the
word ``variety'' in a global application includes a
quasi-projectivity convention, or on the additional sentence
in Corollary 4.7 about preserving a prescribed big open.
\end{proof}
